\documentclass[10pt,a4paper]{amsart}
\usepackage[margin=19mm]{geometry}
\usepackage{amsmath,amssymb,mathtools}
\usepackage[scr=rsfs,cal=euler]{mathalfa}
\usepackage{xfrac}
\usepackage[unicode,hidelinks]{hyperref}
\newcommand{\ie}{i.e.\ }
\newcommand{\cf}{cf.\ }
\newcommand{\resp}{respectively\ }
\newcommand{\Subsection}{\subsection}
\providecommand{\nobreakdash}{\nobreak\hbox{-}}
\setlength{\emergencystretch}{2em}
\multlinegap0pt
\let\ra\rightarrow
\DeclareMathOperator{\supp}{supp}


\newcommand{\RR}{\mathbb R}
\newcommand{\NN}{{\mathbb N}}
\renewcommand{\Im}{\mathop{\mathrm{Im}}}

\newtheorem{theorem}{Theorem}
\newtheorem{lemma}{Lemma}
\newtheorem{proposition}{Proposition}
\newtheorem{corollary}{Corollary}

\theoremstyle{definition}
\newtheorem{remark}{Remark}
\newtheorem{definition}{Definition}

\title{A holographic global uniqueness in passive imaging: Theorem 3 and necessary local core}
\author{Roman G. Novikov}
\date{Source: J.\'E.P.--M. 12 (2025), 1069--1081; DOI 10.5802/jep.306}
\begin{document}\maketitle
\noindent\textit{Editorial scope.} The counted unit is original Theorem 3. Theorems 1--2, Corollary 1 and Propositions 1--2 below are uncounted support. Original equation, theorem, section and citation numbers are retained. The source\textquotesingle s discussions of other measurement surfaces and numerical prospects are omitted.

\section{Introduction}\label{sec:1}

We consider the Helmholtz equation
\begin{equation}
\label{eq:1.1}
-\Delta\psi(x) = \kappa^2 \psi(x),\quad x\in \mathcal{U},\ \kappa>0,
\end{equation}
where $\Delta$ is the Laplacian in $x$, and $\mathcal U$ is an exterior region in $\RR^3$
that is $\mathcal U$ is an open unbounded connected set in $\RR^3$ consisting of all points in exterior of a closed bounded regular surface $S$ (as in \cite{A}, \cite{W}).

For equation \eqref{eq:1.1} we consider the radiation solutions $\psi$ that is the solutions satisfying the Sommerfeld's radiation condition
\begin{equation}\label{eq:1.2}
|x|\Bigl(\frac{\partial}{\partial|x|}-i\kappa\Bigr)\psi(x)\to 0\quad \mbox{as} \ |x|\ra +\infty,
\end{equation}
uniformly in $x/|x|$. We assume that $\psi \in C^2$ in the closure of $\mathcal{U}$.

Let
\begin{equation}\label{eq:1.3}
L=L_{x_0, \theta}=\{x\in\RR^3:\ x=x(s)=x_0+s\theta,\ 0<s <+\infty \}, \quad x_0\in\RR^3, \ \theta\in\mathbb S^2,
\end{equation}
where $\mathbb S^2$ is the unit sphere in $\RR^3$.

The Gelfand-Krein-Levitan problem in question (in its fixed energy version in dimension $d=3$) consists in determining the potential $v$ in the Schrödinger equation
\begin{equation}\label{eq:1.4}
-\Delta\psi(x) +v(x)\psi(x)= \kappa^2 \psi(x) +\delta(x-y),\ x,y\in \RR^3,\ \kappa>0,
\end{equation}
from the imaginary part of the outgoing Green function $R_v^+(x, y, \kappa)$ for one $\kappa$ and all $x$, $y$ on some part of the boundary of a domain containing the support of $v$.
Here, $\delta$ is the Dirac delta function.

\setcounter{equation}{5}
Note that
\begin{equation}\label{eq:1.7}
R_0^+(x, y, \kappa)=\frac{e^{i\kappa|x-y|}}{4\pi|x-y|},\quad x,y\in \RR^3,
\end{equation}
where $R_0^+$ is the outgoing Green function for equation \eqref{eq:1.4} with $v \equiv 0$.

Let
\begin{equation}\label{eq:1.8}
R_{v,sc}^+(x, y, \kappa)=R_v^+(x, y, \kappa)-R_0^+(x, y, \kappa),\quad x,y\in \RR^3.
\end{equation}

Suppose that
\begin{equation}
\label{eq:1.5}
\supp v \subset \RR^3\setminus \overline {\mathcal{U}},
\end{equation}
where $\overline {\mathcal{U}}$ is the closure of $\mathcal{U}$.
Then, in view of the definitions of $R_{v,sc}^+$, $R_v^+$, and $R_0^+$, the function $\psi=R_{v,sc}^+(x, y, \kappa)$ is a radiation solution of equation \eqref{eq:1.1}
for each $y\in \RR^3$.

\section{Main results}\label{sec:2}

In this work our key result is as follows.

\begin{theorem}\label{th:1}
Let $\psi$ be a radiation solution of equation \eqref{eq:1.1} as in \eqref{eq:1.2}.
Let $L$ be a ray as in \eqref{eq:1.3} such that $L \subset \mathcal{U}$, where $\mathcal{U}$ is the region in \eqref{eq:1.1}. Then $\psi$ on $L$ is uniquely determined by $\Im \psi$ on $\Lambda$,
where $\Lambda$ is an arbitrary non-empty open interval of~$L$.
\end{theorem}

As a corollary, we also get, in particular, the following result.

\begin{theorem}\label{th:2}
Let $\psi$ be a radiation solution of equation \eqref{eq:1.1} as in \eqref{eq:1.2}.
Let $X$ be a two-dimensional plane in $\RR^3$ such that $X\subset \mathcal{U}$.
Then $\psi$ on $X$ is uniquely determined by $\Im \psi$ on $\mathcal{D}$,
where $\mathcal{D}$ is an arbitrary non-empty open domain of $X$.
\end{theorem}

Theorem \ref{th:1} is proved in Section \ref{sec:4} using the Atkinson-Wilcox expansion of \cite{A}, \cite{W} for the radiation solutions $\psi$ of equation \eqref{eq:1.1}, and a modified version of holographic techniques of \cite{N1}, \cite{N5}.
In particular, in this proof we use Proposition~\ref{prop:1} of Section~\ref{sec:3}, which yields a two-point approximation for $\psi$
in terms of $\Im \psi$. This two-point approximation is also of independent interest.

Note that $\psi$ and $\Im\psi$ are real-analytic on $\mathcal{U}$, and, therefore, on $L$ in Theorem \ref{th:1} and on $X$ in Theorem \ref{th:2}.
Because of this analyticity, Theorem \ref{th:1} reduces to the case when $\Lambda=L$ and Theorem \ref{th:2} reduces to the case when $\mathcal{D}=X$.

Using this reduction, Theorem \ref{th:2} is proved as follows.
We assume that $\mathcal{D}=X$.
Then to determine $\psi$ at an arbitrary $x\in X$, we consider a ray $L=L_{x_0, \theta}\subset X$ such that $x\in L$ and use Theorem \ref{th:1} for this $L$.

\begin{corollary}\label{cor:1}
Under the assumptions of Theorem \ref{th:2}, the imaginary part $\Im\psi$ on~$\mathcal{D}$, uniquely determines $\psi$
in the entire region $\mathcal{U}$.
\end{corollary}

Corollary \ref{cor:1} follows from Theorem \ref{th:2}, formula \eqref{eq:3.8a} recalled in Section~\ref{subsec:3.3}, and analyticity of $\psi$ in $\mathcal{U}$.

\begin{theorem}\label{th:3}
Let $v \in L^{\infty}(\RR^3)$, $\supp v \subset \Omega$, and $X \subset \RR^3 \setminus \overline {\Omega}$, where $\Omega$ is an open bounded connected domain in $\RR^3$, $\overline {\Omega}$ is the closure of $\Omega$,
and $X$ is a two-dimensional plane. Let property \eqref{eq:3.11} hold for fixed $\kappa>0$ and $R_v^+(x, y, \kappa)$ be the outgoing Green function for equation \eqref{eq:1.4}.
Then $v$ is uniquely determined by $\Im R_v^+(\cdot, \cdot, \kappa)$ on $\mathcal{D} \times \mathcal{D}$,
where $\mathcal{D}$ is an arbitrary non-empty open domain of $X$.
\end{theorem}

In Theorem \ref{th:3} we do not assume that $v$ is real-valued,
but we assume that property~\eqref{eq:3.11} formulated in Section \ref{subsec:3.4} holds.

Theorem \ref{th:3} is proved in Section \ref{sec:5}.
In particular, in this proof we use Proposition~\ref{prop:2} of Section \ref{sec:5},
which stays that $v$ is uniquely determined by $R_v^+(\cdot, \cdot, \kappa)$ on $X\times X$ for fixed~$\kappa$.
To our knowledge, Proposition \ref{prop:2} is slightly different from the results reported in the literature.
Therefore, just in case, for completeness of presentation this proposition is also proved in Section \ref{sec:5}.

\setcounter{equation}{10}
\section{Preliminaries}\label{sec:3}

\subsection{The Atkinson-Wilcox expansion}

Let
\begin{equation}
\label{eq:3.1}
B_r=\{x\in\RR^3:\ |x|<r\},\quad r>0.
\end{equation}
If $\psi$ is a radiation solution of equation \eqref{eq:1.1} and $\RR^3\setminus B_r \subset \mathcal{U}$, then
the following Atkinson-Wilcox expansion holds:
\begin{equation}
\label{eq:3.2}
\psi(x)=\frac{e^{i\kappa|x|}}{|x|}\sum\limits_{j=1}^{\infty}\frac{f_j(\theta)}{|x|^{j-1}} \quad \text{for}\ x\in \RR^3\setminus B_r,\ \theta=\frac{x}{|x|},
\end{equation}
where the series converges absolutely and uniformly; see \cite{A}, \cite{W}.

\subsection{A two-point approximation for $\psi$}\label{subsec:3.2}

Let
\begin{equation}
\label{eq:3.3}
I(x)=|x|\Im\psi(x),\quad x\in\mathcal{U},
\end{equation}
where $\psi$ is a radiation solution of equation \eqref{eq:1.1}.

We have that
\begin{equation}
\label{eq:3.4}
2iI(x)=e^{i\kappa|x|}f_1(\theta)-e^{-i\kappa|x|} \overline {f_1(\theta)}+O(|x|^{-1})\quad \text{as} \ |x|\ra +\infty,
\end{equation}
uniformly in $\theta=x/|x|$, where $f_1$ is the leading coefficient in \eqref{eq:3.2}.

\begin{proposition}\label{prop:1}
Let $\psi$ be a radiation solution of equation \eqref{eq:1.1}. Then
\begin{equation}\label{eq:3.5}
\begin{aligned}
&f_1(\theta)= \frac{1}{\sin(\kappa\tau)}\bigl(-e^{-i\kappa|y|}I(x)+e^{-i\kappa|x|}I(y)+O(|x|^{-1})\bigr),\\
&x,y\in L_{x_0, \theta},\quad x_0=0,\ y=x+\tau\theta, \quad\theta\in\mathbb S^2,\quad \tau>0,
\end{aligned}
\end{equation}
uniformly in $\theta$, where $f_1$ is the coefficient in \eqref{eq:3.4}, $I$ is defined by \eqref{eq:3.3}, $L_{x_0, \theta}$ is the ray defined by \eqref{eq:1.3},
and $\sin(\kappa\tau)\ne 0$ for fixed $\tau$.
\end{proposition}

Formula \eqref{eq:3.5} is a two-point approximation for $f_1$ and together with \eqref{eq:3.2} also gives a two-point approximation for $\psi$
in terms of $I$.
For phaseless inverse scattering and holography formulas of such a type go back to \cite{N1} (see also \cite{NS2}, \cite{NSh}, \cite{N5}).


\setcounter{remark}{3}
\begin{remark}\label{rem:4}
If an arbitrary function $I$ on $L_{0, \theta}$ satisfies \eqref{eq:3.4},
then formula \eqref{eq:3.5} holds, for fixed $\theta\in\mathbb S^2$.
\end{remark}

We obtain \eqref{eq:3.5} from the system of equations for $f_1$ and $\overline {f_1}$:
\begin{equation}\label{eq:3.6}
\begin{aligned}
&e^{i\kappa|x|}f_1(\theta)-e^{-i\kappa|x|} \overline {f_1(\theta)}=2iI(x)+O(|x|^{-1}),\\
&e^{i\kappa|y|}f_1(\theta)-e^{-i\kappa|y|} \overline {f_1(\theta)}=2iI(y)+O(|y|^{-1}),
\end{aligned}
\end{equation}
where $x$, $y$ are as in \eqref{eq:3.5}. In particular, we use that $|y|=|x|+\tau$ and that
\begin{equation}
\label{eq:3.7}
D=2i \sin(\kappa\tau),
\end{equation}
where $D$ is the determinant of system \eqref{eq:3.6}.

In turn, \eqref{eq:3.6} follows from \eqref{eq:3.4}.

\subsection{A Green type formula}\label{subsec:3.3}

The following formula holds:
\begin{align}
\psi(x)&= 2\int_{X}\frac{\partial G^+(x-y,\kappa)}{\partial \nu_{y}}\psi(y)dy,\quad x\in V_X, \label{eq:3.8a}\\
G^+(x,\kappa)&= -\frac{e^{i\kappa|x|}}{4\pi|x|},\quad x\in \RR^3, \label{eq:3.8b}
\end{align}
where $\psi$ is a radiation solution of equation \eqref{eq:1.1}, $X$ and $V_X$ are plane and open half-space in $\mathcal{U}$,
where $X$ is the boundary of $V_X$, $\nu$ is the outward normal to $X$ relative to~$V_X$; see, for example, \cite[Eq.\,(5.84)]{BR}.

Recall that $R_0^+(x, y, \kappa)=-G^+(x-y,\kappa)$, where $R_0^+$ is the outgoing Green function for equation \eqref{eq:1.4} with $v \equiv 0$.

For completeness of presentation note that formula \eqref{eq:3.8a} follows from the formula
\begin{equation}\label{eq:3.9a}
\begin{aligned}
\psi(x)&= \int_{X}\frac{\partial G^+_X(x,y,\kappa)}{\partial \nu_{y}}\psi(y)dy,\quad x\in V_X,\\
G^+_X(x,y,\kappa)&= G^+(x-y,\kappa)-G^+(x-y^*,\kappa),\quad x,y\in V_X \cup X,
\end{aligned}
\end{equation}
where $y^*$ is symmetric to $y$ with respect to $X$.
The point is that $G^+_X$ is the Green function for the Helmholtz operator $\Delta+\kappa^2$ in $V_X$
with Dirichlet boundary condition on $X$ and Sommerfeld radiation condition at infinity.

\subsection{Some facts of direct scattering}\label{subsec:3.4}

We consider equation \eqref{eq:1.4} assuming for simplicity that
\begin{equation}\label{eq:3.9}
\begin{aligned}
&v \in L^{\infty}(\Omega),\quad v \equiv 0\ \mbox{on}\ \RR^3\setminus \overline {\Omega}\\
&\Omega\ \mbox{is\ an\ open\ bounded\ connected\ domain\ in}\ \RR^3.
\end{aligned}
\end{equation}

The outgoing Green function $R_v^+$ for equation \eqref{eq:1.4} satisfies the integral equation
\begin{equation}
\label{eq:3.10}
R_v^+(x, y, \kappa)=-G^+(x-y,\kappa) +\int_{\Omega}G^+(x-z,\kappa)v(z)R_v^+(z, y, \kappa)dz,
\end{equation}
where $x,y \in \RR^3$, $G^+$ is given by \eqref{eq:3.8b}.

Actually, in addition to \eqref{eq:3.9}, we assume that, for fixed $\kappa>0$,
\begin{equation}
\label{eq:3.11}
\mbox{equation}\ \eqref{eq:3.10}\ \mbox{is\ uniquely\ solvable\ for}\ R_v^+(\cdot , y, \kappa) \in L^2(\Omega).
\end{equation}
In particular, it is known that if $v$ satisfies \eqref{eq:3.9} and is real-valued (or $\Im v \leq 0$), then~\eqref{eq:3.11} is fulfilled automatically;
see, for example, \cite{CK}.

We also consider the scattering wave functions $\psi^+$ for the homogeneous equation~\eqref{eq:1.4} (\ie without $\delta$):
\begin{equation}
\label{eq:3.12}
\psi^+=\psi^+(x,\theta,\kappa)=e^{i \kappa\theta x}+\psi^+_{sc}(x,\theta,\kappa),\quad x \in \RR^3,\ \theta \in \mathbb S^2,
\end{equation}
where $\psi^+_{sc}$ satisfies the radiation condition \eqref{eq:1.2} at fixed $\theta$.

The following formulas hold:
\begin{align}
\label{eq:3.13}
R_v^+(x, y, \kappa)&=R_v^+(y, x, \kappa),\quad x, y\in \RR^3;
\\
\label{eq:3.14}
R_v^+(x, y, \kappa)&=\frac{e^{i\kappa|x|}}{4\pi|x|}\,\psi^+\Bigl(y, -\frac{x}{|x|},\kappa\Bigr)+ O\Bigl(\frac{1}{|x|^2}\Bigr) \quad \mbox{as} \ |x|\ra +\infty\text{ at fixed $y$;}\\
\label{eq:3.15}
\psi^+_{sc}(x,\theta,\kappa)&=\frac{e^{i\kappa|x|}}{|x|}\,A\Bigl(\theta, \frac{x}{|x|},\kappa\Bigr)+ O\Bigl(\frac{1}{|x|^2}\Bigr) \quad \mbox{as} \ |x|\ra +\infty\text{ at fixed $\theta$,}
\end{align}
where $A$ arising in \eqref{eq:3.15} is the scattering amplitude for the homogeneous equation~\eqref{eq:1.4}
and is defined on $\mathbb S^2 \times \mathbb S^2$ at fixed $\kappa$.

In view of \eqref{eq:1.7}, \eqref{eq:1.8}, \eqref{eq:3.12}--\eqref{eq:3.14}, we also have that
\begin{align}
\label{eq:3.16}
R_{v,sc}^+(x, y, \kappa)&=R_{v,sc}^+(y, x, \kappa),\ x, y\in \RR^3;
\\
\label{eq:3.17}
R_{v,sc}^+(x, y, \kappa)&=\frac{e^{i\kappa|x|}}{4\pi|x|}\,\psi_{sc}^+\Bigl(y, -\frac{x}{|x|},\kappa\Bigr)+ O\Bigl(\frac{1}{|x|^2}\Bigr) \ \mbox{as} \ |x|\ra +\infty\text{ at fixed $y$.}
\end{align}
In connection with aforementioned facts concerning $R_v^+$ and $\psi^+$ see, e.g., \cite[Chap.\,IV, \S1]{FM}.

\begin{remark}\label{rem:5}
It is well known that, under assumptions \eqref{eq:3.9}, \eqref{eq:3.11}, the scattering amplitude $A=A(\cdot, \cdot, \kappa)$ is real analytic on $\mathbb S^2 \times \mathbb S^2$.
\end{remark}

\section{Proof of Theorem \ref{th:1}}\label{sec:4}

\subsection{Case $L\subseteq L_{0, \theta}$}\label{subsec:4.1}
First, we give the proof for the case when $L\subseteq L_{0, \theta}$.
In this case it is essentially sufficient to prove that $\Im\psi$ on $L$ uniquely determines $f_j(\theta)$ in~\eqref{eq:3.2} for all $j$.
Such a determination is presented below in this subsection.
The rest follows from the convergence of the series in \eqref{eq:3.2} and analyticity of $\psi$ and $\Im\psi$ on~$L$.

The determination of $f_1$ follows from \eqref{eq:3.5}.
Suppose that $f_1,\dots,f_n$ are determined, then the determination of $f_{n+1}$ is as follows.
Let
\begin{align}
\label{eq:4.1}
\psi_n(x)&=\frac{e^{i\kappa|x|}}{|x|}\sum\limits_{j=1}^{n}\frac{f_j(\theta)}{|x|^{j-1}}, \quad \text{where}\ \theta=\frac{x}{|x|},
\\
\label{eq:4.2}
I_n(x)&=|x|\Im\psi_n,
\\
\label{eq:4.3}
J_n(x)&=|x|^n(I(x)-I_n(x)),
\end{align}
where $x$ is as in \eqref{eq:3.2}, $I(x)$ is defined by \eqref{eq:3.3}.

We have that
\begin{align}
\label{eq:4.4}
2iI(x)&=2iI_n(x)+ \frac{e^{i\kappa|x|}}{|x|^n}{f_{n+1}(\theta)}- \frac{e^{-i\kappa|x|}}{|x|^n}\,\overline{f_{n+1}(\theta)}+O(|x|^{-n-1}),
\\
\label{eq:4.5}
2iJ_n(x)&=e^{i\kappa|x|}{f_{n+1}(\theta)} - e^{-i\kappa|x|}\overline{f_{n+1}(\theta)}+O(|x|^{-1}),
\end{align}
as $|x|\to +\infty$, where $I$ is defined by \eqref{eq:3.3}.

Due to \eqref{eq:4.5} and Remark \ref{rem:4}, we get
\begin{equation}\label{eq:4.6}
\begin{aligned}
&f_{n+1}(\theta)= \frac{1}{\sin(\kappa\tau)}\bigl(-e^{-i\kappa|y|}J_n(x)+e^{-i\kappa|x|}J_n(y)+O(|x|^{-1})\bigr),\\
&x,y\in L_{x_0, \theta},\quad x_0=0,\quad y=x+\tau\theta, \quad \theta\in\mathbb S^2,\quad \tau>0,
\end{aligned}
\end{equation}
assuming that $\sin(\kappa\tau)\ne 0$ for fixed $\tau$
(where the parameter $\tau$ can be always fixed in such a way for fixed $\kappa>0$).

Formulas \eqref{eq:3.3}, \eqref{eq:4.1}--\eqref{eq:4.3} and \eqref{eq:4.6} determine $f_{n+1}$, give the step of induction for finding all $f_j$, and complete the proof of Theorem \ref{th:1} for the case $L\subseteq L_{0, \theta}$.

\subsection{General case}
The general case reduces to the case of Section \ref{subsec:4.1} by the change of variables
\begin{equation}
x'=x-q \label{eq:4.7}\quad
\text{for some fixed $q\in\RR^3$ such that $L\subseteq L_{q, \theta}$.}
\end{equation}
In the new variables $x'\in\mathcal{U}'=\mathcal{U}-q$, we have that:
\begin{equation}
\label{eq:4.8}
\psi= \frac{e^{i\kappa|x'|}}{|x'|}\sum\limits_{j=1}^{\infty}\frac{f_j'(\theta)}{|x'|^{j-1}} \quad \text{for}\ x'\in \RR^3\setminus B_{r'}, \ \theta=\frac{x'}{|x'|},
\end{equation}
for some new $f_j'$,
where $r'$ is such that $\RR^3\setminus B_{r'} \subset \mathcal{U}'$;
\begin{equation}
\label{eq:4.9}
L\subseteq L_{q, \theta}=L_{0, \theta}.
\end{equation}

In addition, the series in \eqref{eq:4.8} converges absolutely and uniformly.

In view of \eqref{eq:4.8}, \eqref{eq:4.9},
we complete the proof of Theorem \ref{th:1} by repeating the proof of Section \ref{subsec:4.1}.

\section{Proof of Theorem \ref{th:3}}\label{sec:5}


Under our assumptions on $v$, $\Omega$, $X$, and $\kappa$, we have, in particular, that
\begin{align}
&R_{v,sc}^+(\cdot, y, \kappa)\text{ is real-analytic on}\ X\text{ for fixed}\ y \in X, \label{eq:5.1a} \\
&R_{v,sc}^+(x, \cdot, \kappa)\text{ is real-analytic on}\ X\text{ for fixed}\ x \in X, \label{eq:5.1b}
\end{align}
where $R_{v,sc}^+$ is defined by \eqref{eq:1.7}, \eqref{eq:1.8}.
Here, we use that $R_{v,sc}^+(x, y, \kappa)$ satisfies the homogeneous equation \eqref{eq:1.4} and, therefore, is real-analytic outside of $\supp v$,
and that symmetry \eqref{eq:3.16} holds.

In view of \eqref{eq:1.7}, \eqref{eq:1.8}, \eqref{eq:5.1a}, \eqref{eq:5.1b}, Theorem \ref{th:3} reduces to the case when $\mathcal{D}=X$.
In turn, Theorem \ref{th:3} with $\mathcal{D}=X$ follows from formulas \eqref{eq:1.7}, \eqref{eq:1.8} and from Lemma \ref{lem:1} and Proposition~\ref{prop:2} given below.

\begin{lemma}\label{lem:1}
Under the conditions of Theorem \ref{th:3}, $\Im R_{v,sc}^+(\cdot, y, \kappa)$ on $X$ uniquely determines $R_{v,sc}^+(\cdot, y, \kappa)$ on $X$,
where $y \in \RR^3$, and $\Im R_{v,sc}^+(\cdot, \cdot, \kappa)$ on $X\times X$ uniquely determines $R_{v,sc}^+(\cdot, \cdot, \kappa)$ on $X\times X$.
\end{lemma}

Lemma \ref{lem:1} follows from Theorem \ref{th:2} with $\mathcal U$ such that $\overline {\Omega} \subset \RR^3\setminus \overline {\mathcal{U}}$, $X \subset \mathcal U$, and
the property that $\psi=R_{v,sc}^+(x, y, \kappa)$ is a radiation solution of equation \eqref{eq:1.1} for each $y\in \RR^3$.

\begin{proposition}\label{prop:2}
Under the conditions of Theorem \ref{th:3}, $R_{v,sc}^+(\cdot, \cdot, \kappa)$ on $X\times X$ uniquely determines $v$ (for fixed $\kappa$).
\end{proposition}

Proposition \ref{prop:2} is proved as follows (for example).
First,
\begin{equation}\label{eq:5.2}
\begin{aligned}
&\psi=R_{v,sc}^+(\cdot, x', \kappa)\text{ on}\ X\text{ uniquely determines}\ \psi=R_{v,sc}^+(\cdot, x', \kappa)\text{ on}\ V_X\\
&\mbox{\rm via formula}\ \eqref{eq:3.8a}, \text{ for each fixed }\ x' \in X.
\end{aligned}
\end{equation}

Second,
\begin{equation}\label{eq:5.3}
\begin{aligned}
&R_{v,sc}^+(\cdot, x', \kappa)\text{ on}\ V_X \cup X\text{ uniquely determines}\ \psi_{sc}^+(x',\theta,\kappa)\text{ for}\ \theta \in \Theta_X^+\\
&\mbox{\rm via formula}\ \eqref{eq:3.17}, \text{ for each fixed}\ x' \in X,
\end{aligned}
\end{equation}
where
\begin{equation}
\label{eq:5.4}
\Theta_X^{\pm}=\{\theta \in\mathbb S^2:\ \pm\theta\nu \geq 0 \},
\end{equation}
where $\nu$ is the outward normal to $X$ relative to $V_X$ considered as the interior of $X$.

Third,
\begin{equation}
\begin{aligned}
&\psi=\psi_{sc}^+(\cdot, \theta,\kappa)\text{ on}\ X\text{ uniquely determines}\ \psi=\psi_{sc}^+(\cdot, \theta,\kappa)\text{ on}\ V_X \label{eq:5.5}\\
&\text{via formula \eqref{eq:3.8a}}.
\end{aligned}
\end{equation}

Fourth,\vspace*{-3pt}
\begin{equation}\label{eq:5.6}
\begin{aligned}
&\psi_{sc}^+(\cdot, \theta,\kappa)\text{ on}\ V_X \cup X\text{ uniquely determines}\ A(\theta, \theta', \kappa)\text{ for}\ \theta' \in \Theta_X^-\\
&\text{via formula \eqref{eq:3.15},  for each fixed $\theta \in\mathbb S^2$,}
\end{aligned}
\end{equation}
where $\Theta_X^-$ is defined in \eqref{eq:5.4}.

Therefore, we get that\vspace*{-3pt}
\begin{equation}\label{eq:5.7}
\begin{aligned}
&R_{v,sc}^+(\cdot, \cdot, \kappa)\text{ on}\ X\times X\text{ uniquely determines}\ A(\cdot, \cdot, \kappa)\text{ on}\ \Theta_X^+ \times \Theta_X^-\\
&\text{via \eqref{eq:5.2}, \eqref{eq:5.3}, \eqref{eq:5.5}, \eqref{eq:5.6}.}
\end{aligned}
\end{equation}
In addition,\vspace*{-3pt}
\begin{equation}\label{eq:5.8}
\begin{aligned}
&A(\cdot, \cdot, \kappa)\text{ on}\ \Theta_X^+ \times \Theta_X^-\text{ uniquely determines}\ A(\cdot, \cdot, \kappa)\text{ on}\ \mathbb S^2 \times \mathbb S^2 \\
&\text{by analyticity},
\end{aligned}
\end{equation}
in view of Remark \ref{rem:5}.

Finally, Proposition \ref{prop:2} follows from \eqref{eq:5.7}, \eqref{eq:5.8} and the result of \cite{N1988} (see also \cite{N1994}, \cite{HH}, \cite{E}) that,
under assumptions \eqref{eq:3.9}, \eqref{eq:3.11}, the scattering amplitude $A$ at fixed $\kappa$ uniquely determines $v$.
In this result the assumption that $v$ is real-valued or $\Im v \leq 0$ is not essential and can be replaced by assumption~\eqref{eq:3.11}.

This completes the proofs of Proposition \ref{prop:2} and Theorem \ref{th:3}.

\begin{thebibliography}{Zim84}
\bibitem[3]{A} F. V. Atkinson, ``On Sommerfeld\textquotesingle s radiation condition'', \emph{Philos. Math.} (7) \textbf{40} (1949), 645--651.
\bibitem[5]{BR} V. Burov and O. Rumyantseva, \emph{Inverse wave problems of acoustic tomography. Part 2, Inverse problems of acoustic scattering}, LENAND, Moscow, 2019, in Russian.
\bibitem[7]{CK} D. Colton and R. Kress, \emph{Inverse acoustic and electromagnetic scattering theory}, fourth ed., Applied Math. Sciences 93, Springer, Cham, 2019.
\bibitem[8]{E} G. Eskin, \emph{Lectures on linear partial differential equations}, Graduate Studies in Math. 123, American Mathematical Society, Providence, RI, 2011.
\bibitem[9]{FM} L. D. Faddeev and S. P. Merkuriev, \emph{Quantum scattering theory for several particle systems}, Mathematical Physics and Applied Math. 11, Kluwer Academic Publishers Group, Dordrecht, 1993, translated from the 1985 Russian original.
\bibitem[12]{HH} P. H\"ahner and T. Hohage, ``New stability estimates for the inverse acoustic inhomogeneous medium problem and applications'', \emph{SIAM J. Math. Anal.} \textbf{33} (2001), no. 3, 670--685.
\bibitem[19]{N1988} R. G. Novikov, ``A multidimensional inverse spectral problem for the equation $-\Delta\psi+(v(x)-Eu(x))\psi=0$'', \emph{Funktsional. Anal. i Prilozhen.} \textbf{22} (1988), no. 4, 11--22.
\bibitem[20]{N1994} R. G. Novikov, ``The inverse scattering problem at fixed energy for the three-dimensional Schr\"odinger equation with an exponentially decreasing potential'', \emph{Comm. Math. Phys.} \textbf{161} (1994), no. 3, 569--595.
\bibitem[21]{N1} R. G. Novikov, ``Inverse scattering without phase information'', in \emph{S\'eminaire Laurent Schwartz---\'Equations aux d\'eriv\'ees partielles et applications. Ann\'ee 2014--2015}, Editions de l\textquotesingle\'Ecole Polytechnique, Palaiseau, 2016, XVI.1--XVI.13.
\bibitem[23]{N5} R. G. Novikov, ``A holographic uniqueness theorem'', \emph{Proc. Steklov Inst. Math.} \textbf{325} (2024), 218--223.
\bibitem[24]{NSh} R. G. Novikov and B. L. Sharma, ``Phase recovery from phaseless scattering data for discrete Schr\"odinger operators'', \emph{Inverse Problems} \textbf{39} (2023), no. 12, article 125006 (12 pages).
\bibitem[25]{NS2} R. G. Novikov and V. N. Sivkin, ``Fixed-distance multipoint formulas for the scattering amplitude from phaseless measurements'', \emph{Inverse Problems} \textbf{38} (2022), no. 2, article 025012 (22 pages).
\bibitem[29]{W} C. H. Wilcox, ``A generalization of theorems of Rellich and Atkinson'', \emph{Proc. Amer. Math. Soc.} \textbf{7} (1956), 271--276.
\end{thebibliography}
\end{document}
